\section{Conclusion}

We presented the first controlled comparison of distillation objective types for Transformer-to-SSM conversion, revealing that token-level alignment produces representations 5$\times$ more stable across sequence lengths than matrix-level alignment.

Our unified Phi-Mamba framework enabled fair comparison between MOHAWK (matrix-level) and CAB (token-level) objectives. Measuring hidden state drift across sequence lengths 512--2048, we found:

\begin{itemize}
    \item CAB drift slope: 0.00090506
    \item MOHAWK drift slope: 0.00452495
    \item Ratio: 5$\times$
\end{itemize}

This stability gap reflects a mechanistic difference: token-level supervision is position-agnostic, while matrix-level supervision encodes position-position relationships that change with sequence length.

\subsection{Future Work}

Three directions emerge from this work:

\textbf{Full training verification}: Our PoC validates methodology but cannot verify downstream F1 advantage. Full 1.5B token training per condition would test whether the 5$\times$ stability gap translates to task performance at extrapolated lengths.

\textbf{RoPE-based teachers}: Llama, Mistral, and other RoPE-based models can process 16K--32K sequences natively. Extending our comparison to these teachers would enable true length extrapolation experiments.

\textbf{Hybrid objectives}: Both MOHAWK and CAB objectives can coexist in our unified framework. Combining matrix-level supervision at training lengths with token-level supervision for extrapolation may capture benefits of both approaches.

The 5$\times$ stability gap we identify provides a principled starting point for these investigations---a quantified criterion for objective selection in Transformer-to-SSM distillation.
